\section{Conclusion}

This thesis studies how a major online chess tournament reform changed the relationship between player characteristics, decision quality, and competitive outcomes. In September 2025, Chess.com changed Titled Tuesday from a 3+1 increment format with two weekly time slots to a 5+0 no-increment format with one main weekly slot. The reform created a useful empirical setting because it affected a repeated elite competition with stable player identities, objective ratings, complete game outcomes, and reconstructable move histories.

The strongest result is that the new format increased the return to relative skill. A 100-point rating advantage became worth about 1.2 additional percentage points of score share after the reform. This result survives the main robustness checks: excluding early rounds, restricting the window around the reform, using balanced players, adding paired-game fixed effects, and comparing against placebo cutoffs. The win/loss decomposition shows that stronger players won more and lost less, while draw probability changed little.

The second major result is that outcomes became more tightly linked to realized move quality. Accuracy advantages converted more strongly into points under 5+0. This is not interpreted as a causal effect of accuracy because accuracy is measured inside the game. Instead, it describes a steeper chess production function: given the same observed quality advantage over an opponent, the advantage was worth more in the no-increment format.

The move-level and clock analysis explains where this skill premium appears. Late-game and final-ten-ply errors increased under 5+0. Critical equal positions became more costly. Rating favorites converted winning positions more effectively, converted +2 positions faster, let fewer winning positions fall back to equality, and recovered more from -2 positions. Direct clock evidence sharpens the interpretation. Ordinary low-clock exposure became less common because players started with more time, but severe time trouble became more costly, recovery from low-clock states nearly disappeared, and opponent clock weakness became more valuable for conversion and defensive escape. Predetermined defensive skill and complexity tolerance predicted lower post-change errors. These results support a conversion-technology interpretation: 5+0 rewards players who can turn advantages into wins and survive worse positions under a different time-control technology.

Adaptability is heterogeneous. Younger players performed better relative to older players after the reform, although pretrend and placebo evidence require cautious causal language. Online-platform capital also mattered: players with high Chess.com ratings relative to classical ratings gained more, consistent with the value of interface fluency and online practical skill. Old late-slot regulars were harmed by the schedule component of the reform, but mainly through participation and completion margins rather than lower per-game accuracy.

The style analysis adds a data-science layer. Interpretable pre-change style indexes show that historically chaotic profiles did not receive a relative score premium from 5+0. If anything, chaos-creation style is associated with lower post-change score share. This rejects a plausible story that no-increment chess rewards practical chaos. The neural contrastive model shows that stable player-style embeddings can be learned from pre-change move sequences, and the resulting clusters align with meaningful chess profiles. All five neural style clusters show raw post-change accuracy improvements, so the neural result is not a null result: it shows that the move-quality gains from the new format were broad across learned styles, even though chaotic style did not become relatively more rewarded in score terms.

Several hypotheses were not supported. The data do not show robust effects for GDP, White advantage, broad female-status interactions, country local-time exposure, opening preparation, or simplification strategy. They also do not support uniform complex-position deterioration, repeated-opponent familiarity, or post-blunder error cascades as main channels. These null results are useful because they narrow the interpretation. The reform did not simply make chess more random, nor did it mostly reward openings or first-move initiative. It appears to have increased the value of skillful conversion and online adaptation while also changing who could participate consistently.

The main limitation is that the reform bundled the time-control change with the removal of the late slot and changes in the event's broader competitive cycle. The analysis separates conditional game performance from participation and completion, but the institutional shock is still multi-dimensional. A second measurement caveat is that clock annotations measure recorded remaining time, not psychological pressure directly. Premoves, interface latency, and online execution speed can affect the mapping between clock change and cognitive thinking time.

Future work could extend the thesis in three directions. First, the clock models could be combined with richer behavioral data on premoves, mouse speed, and interface latency to separate thinking time from execution time. Second, the neural style model could be extended from descriptive cluster comparisons to a full econometric treatment-heterogeneity design by estimating post-change style-cluster interactions with player and event fixed effects. Third, the tournament-selection part could be expanded into a structural participation model that separates time-zone costs, prize incentives, expected field strength, and player availability.

The overall empirical conclusion is that the Titled Tuesday reform changed both performance and selection. Conditional on playing, skill mattered more. In the move and clock data, that skill premium is visible in conversion, defensive recovery, late-game mistakes, critical positions, severe low-clock states, and the disappearance of increment insurance. In the tournament data, the single-slot design shifted participation and completion. The format change therefore reweighted elite online chess toward players who combine objective strength, online platform capital, and stable conversion under no-increment constraints.
